\documentclass[11pt]{article}
\usepackage[utf8]{inputenc}
\usepackage{geometry}
\usepackage{hyperref}
\usepackage{enumitem}
\geometry{a4paper, margin=1in}

\title{\textbf{DAO Governance x Atala PRISM\\--- Project Close-Out Report ---}}
\author{MuesliSwap Team}
\date{October 2025}

\begin{document}

\maketitle

\section*{Official Close-Out Summary}
\begin{itemize}[leftmargin=1cm]
    \item \textbf{Name of project and Project URL on Fund:}\\
    DAO Governance x Atala PRISM\\
    {\scriptsize\url{https://projectcatalyst.io/funds/10/atala-prism-launch-ecosystem/dao-governance-x-atala-prism-by-muesliswap}}

    \item \textbf{Project ID:} 1000059

    \item \textbf{Project managed by:} MuesliSwap Team

    \item \textbf{Project started in:} October 2023

    \item \textbf{Project completed in:} October 2025


    \item \textbf{Challenge KPIs and how the project addressed them:}
    \begin{itemize}
        \item \textbf{Decentralization and Governance:} Delivered open-source smart contracts and infrastructure enabling on-chain treasury management and parameter updates governed by community voting, directly advancing decentralized decision-making on Cardano.
        \item \textbf{Atala PRISM Adoption:} Implemented DID integration within the governance framework, demonstrating practical use of Atala PRISM for identity-secured, tokenless voting systems.
    \end{itemize}

    \item \textbf{Project KPIs and how the project addressed them:}
    \begin{itemize}
        \item \textbf{Atala PRISM DID Integration:} Completed the conceptual design and implemented DID-based authentication for off-chain voting.
        \item \textbf{On-chain Treasury Contract:} Developed, tested, and internally audited a treasury smart contract allowing community-controlled fund management.
        \item \textbf{Protocol Parameter Governance:} Designed and demonstrated integration of parameter update mechanisms for DeFi protocols governed via smart contracts.
        \item \textbf{Open-source Deliverables:} Released all codebases, documentation, and integration guides publicly on GitHub.
        \item \textbf{Community Engagement:} Shared progress and findings through blog posts, Twitter Spaces, and open development discussions.
    \end{itemize}

    \item \textbf{Key achievements:}
    \begin{itemize}
        \item Built a unified governance framework supporting authenticated on-chain voting.
        \item Successfully integrated Atala PRISM DIDs into governance workflows.
        \item Delivered fully functional treasury and parameter update smart contracts.
        \item Published detailed documentation and tutorials for adoption by other Cardano projects.
    \end{itemize}

    \item \textbf{Key learnings:}
    \begin{itemize}
        \item Combining identity-based and token-based governance models enhances flexibility and fairness in DAOs.
        \item Designing secure treasury and governance contracts requires careful handling of transaction conditions and voting verification logic.
        \item Transparent community engagement during development improves both adoption and trust.
    \end{itemize}

    \item \textbf{Next steps for the product or service developed:}
    \begin{itemize}
        \item Expand integration of the governance framework into more Cardano DeFi protocols.
        \item Collaborate with other projects to deploy the treasury module for real-world DAO management.
        \item Conduct an open testing phase with community participation.
    \end{itemize}

    \item \textbf{Final thoughts/comments:}\\
    The DAO Governance × Atala PRISM project delivered a foundational, open-source governance infrastructure for Cardano. By combining decentralized identity with on-chain treasury management and parameter control, this project strengthens the basis for fully decentralized DAOs. We believe the developed tools will play a central role in shaping the next generation of transparent, community-driven governance on Cardano.

    \item \textbf{Links to other relevant project sources or documents:}
    \begin{itemize}
        \item GitHub Repository:\\ \url{https://github.com/MuesliSwapTeam/muesliswap-atala-onchain-governance}
        \item Close-Out Video:\\ \url{https://youtu.be/3uE83l8lkv4}
    \end{itemize}
\end{itemize}

\section{Introduction}
This report summarizes the development of an open-source governance framework integrating Atala PRISM decentralized identities (DIDs) with on-chain voting mechanisms.
The project extends the governance platform previously developed by the MuesliSwap team and adds new components for decentralized treasury management and protocol parameter updates.  
The goal was to explore how DID-based identity can complement or replace token-based voting and to build the necessary smart contract and backend infrastructure to support this functionality.

\section{Project Overview and Milestones}
The project was carried out over several months and divided into five milestones, progressing from conceptual research to the implementation of advanced governance mechanisms.

\begin{enumerate}
    \item \textbf{Milestone 1 – Research:}\\
    This phase focused on the conceptual and analytical foundation of the project.  
    The team sketched how Atala PRISM DIDs could be integrated into the existing off-chain voting system, designed the basic functionality required for a treasury contract and a protocol parameter update interface, and reviewed existing on-chain governance approaches on Cardano.  
    The output was a publicly available document outlining proposed designs, comparisons with existing systems, and reasoning behind selected technical directions.

    \item \textbf{Milestone 2 – Atala PRISM DID Integration:}\\
    In this phase, the team implemented DID support in the governance system.  
    This included a snapshot mechanism for elections, backend infrastructure to handle DID-based identities, and frontend components for user interaction.  
    The resulting system allowed DID-based participation in the voting process.  
    The code was released as open source, accompanied by documentation and a short demonstration video.

    \item \textbf{Milestone 3 – Treasury Smart Contract:}\\
    This milestone introduced a treasury contract for decentralized fund management.  
    The team developed both the on-chain contract and supporting off-chain infrastructure, performed testnet deployments, and conducted an internal security audit.  
    The final audited version was made publicly available on GitHub.

    \item \textbf{Milestone 4 – On-Chain Voting:}\\
    The fourth milestone focused on integrating Atala DID-based authentication with the existing on-chain platform.
    Internal auditing confirmed correct functionality and no critical vulnerabilities, and the system was tested through trial elections.

    \item \textbf{Final Milestone – Protocol Parameter Update Mechanisms:}\\
    The final stage integrated parameter update capabilities into a selected open-source DeFi protocol, namely a simple AMM (automated market making) pool.  
    This required modifications to smart contracts, off-chain tooling, and additional security and functionality tests.  
    Documentation and tutorials were prepared to support future adoption, alongside community outreach through a blog post, post on social media, and a close-out video.  
    Successful example transactions on testnet demonstrated the full governance flow—from DID-based voting to protocol parameter updates.
\end{enumerate}

\section{Conclusion}
The project successfully delivered a modular governance framework combining identity-based voting, treasury management, and on-chain parameter control.  
All milestones were completed as planned, with open-source releases, documentation, and testing reports made publicly available.  
The resulting system provides a practical foundation for future experiments and implementations of decentralized governance within the Cardano ecosystem.

\bigskip

\noindent \textbf{For further information or to contribute, please refer to the links provided in the summary above.}

\end{document}